\chapter{Vulnerability Report}

\section{Overview}
This chapter presents the findings of the vulnerability assessment. Each identified vulnerability is documented with its corresponding severity level, CVE reference, and CVSS score, technical descriptions, and targeted remediation strategies.
\\

The CVSS 3.0 score and Risk Factor for known vulnerabilities were sourced from NIST National Vulnerability Database \cite{nvd}.
\\

To evaluate novel vulnerabilities, the NIST CVSS Calculator was used \cite{cvsscalc}.
\\


The assessment identified \textbf{89 vulnerabilities}:
\\

\begin{center}
\begin{tabular}{lr}
\toprule
\textbf{Severity} & \textbf{Count} \\
\midrule
Critical & 21 \\
High & 45 \\
Medium & 19 \\
Low & 0 \\
Not Rated/Info & 4 \\
\bottomrule
\end{tabular}
\end{center}

\section{Vulnerabilities}

%%%%%%%%%%%%%%%%%%%%%%%%%%%%%%%%%%%%%%%%%%%%
%
% Corrosion
%
%%%%%%%%%%%%%%%%%%%%%%%%%%%%%%%%%%%%%%%%%%%%
\subsection{Multiple Corrosion Weaknesses}

\subsubsection{F-001: Sensitive Information Exposure in Public Web Root}
\begin{tabular}{|p{2.8cm}|p{12.2cm}|}
\hline
\textbf{CWEs} & CWE-200, CWE-219, CWE-538, CWE-552 \\ \hline
\textbf{CVSS 3.0} & 8.6 \\ \hline
\textbf{Risk Factor} & HIGH \\ \hline
\textbf{Description} & The web application stores files containing potentially sensitive operational data, including log files, to-do lists, and developer notes, within the document root directory. These files are directly accessible to external, unauthenticated users via standard HTTP requests. Exposed paths such as \texttt{/tasks/} and \texttt{/blog-post/archives/} allow unauthorized discovery of internal system details. \\ \hline
\textbf{Solution} & \textbf{Restrict Access:} Move sensitive files (\texttt{tasks\_todo.txt}, log files, configuration files) outside the public web root directory (\texttt{wwwroot}, \texttt{public\_html}). \newline \textbf{Access Controls:} Implement HTTP authentication (e.g., \texttt{.htaccess} rules or web server access directives) to restrict access to sensitive endpoints. \newline \textbf{Disable Directory Browsing:} Ensure directory indexing is explicitly disabled across the web server configuration to prevent unauthorized asset discovery. \\ \hline
\end{tabular}

\subsubsection{F-002: Local File Inclusion (LFI) via Path Traversal in randylogs.php}
\begin{tabular}{|p{2.8cm}|p{12.2cm}|}
\hline
\textbf{CWEs} & CWE-22, CWE-98 \\ \hline
\textbf{CVSS 3.0} & 8.6 \\ \hline
\textbf{Risk Factor} & HIGH \\ \hline
\textbf{Description} & The \texttt{randylogs.php} script accepts untrusted user input via the \texttt{file} parameter and passes it directly into PHP file inclusion or file reading functions without restricting execution scope. The script fails to sanitize path parameters, allowing attackers to perform directory traversal and read arbitrary files outside the web root directory, such as \texttt{/var/log/auth.log}. \\ \hline
\textbf{Solution} & \textbf{Fix LFI:} Do not pass user input directly into file-handling functions (\texttt{include}, \texttt{require}, \texttt{file\_get\_contents}, \texttt{readfile}). \newline \textbf{Strict Whitelisting:} Implement a strict whitelist mapping fixed parameter keys to predefined file paths instead of accepting dynamic file paths from users. \\ \hline
\end{tabular}

\subsubsection{F-003: Log Poisoning via Unsanitized SFTP Authentication Input}
\begin{tabular}{|p{2.8cm}|p{12.2cm}|}
\hline
\textbf{CWEs} & CWE-20, CWE-116 \\ \hline
\textbf{CVSS 3.0} & 9.3 \\ \hline
\textbf{Risk Factor} & CRITICAL \\ \hline
\textbf{Description} & The SFTP service fails to sanitize user input supplied during authentication attempts before writing entries to system logs (e.g., \texttt{/var/log/auth.log}). An attacker can inject executable PHP directives into authentication requests. When these poisoned log files are subsequently processed by vulnerable scripts such as \texttt{randylogs.php}, the server interprets the injected payload as executable PHP code, enabling remote code execution. \\ \hline
\textbf{Solution} & \textbf{Input Neutralization:} Sanitize and escape all user-supplied input prior to writing data to system or application log files. \newline \textbf{Restrict Log Access:} Ensure that log files generated by system services are not readable by the web server process user account. \\ \hline
\end{tabular}

\subsubsection{F-004: Insecure Local Permissions on Backup Artifacts}
\begin{tabular}{|p{2.8cm}|p{12.2cm}|}
\hline
\textbf{CWEs} & CWE-530 \\ \hline
\textbf{CVSS 3.0} & 5.3 \\ \hline
\textbf{Risk Factor} & MEDIUM \\ \hline
\textbf{Description} & The system backup file \texttt{/var/backups/user\_backups.zip} is stored with overly permissive file permissions. The file is readable by the low-privileged web server service account (\texttt{www-data}), exposing sensitive internal backup artifacts to any user or process with local read access. \\ \hline
\textbf{Solution} & \textbf{Restrict File Permissions:} Apply strict file permissions (e.g., \texttt{chmod 600}) to backup archives so they are accessible only by the \texttt{root} user or dedicated system administration accounts. \newline \textbf{Secure Backup Path:} Store system archives in protected directories outside shared system locations. \\ \hline
\end{tabular}

\subsubsection{F-005: Weak Password Architecture on Backup Archives}
\begin{tabular}{|p{2.8cm}|p{12.2cm}|}
\hline
\textbf{CWEs} & CWE-521 \\ \hline
\textbf{CVSS 3.0} & 7.8 \\ \hline
\textbf{Risk Factor} & HIGH \\ \hline
\textbf{Description} & The encrypted archive \texttt{/var/backups/user\_backups.zip} relies on a simple, dictionary-based password pattern. This lack of password complexity allows local or remote attackers who retrieve the archive to perform rapid offline password cracking to extract protected contents. \\ \hline
\textbf{Solution} & \textbf{Secure Data Encryption:} Replace weak password patterns and legacy ZIP encryption with strong AES-256 encryption using high-entropy generated keys. \newline \textbf{Password Policy:} Enforce strong length and complexity requirements for all administrative and automated system archive credentials. \\ \hline
\end{tabular}

\subsubsection{F-006: Local Privilege Escalation via SUID Binary Search Path Hijacking}
\begin{tabular}{|p{2.8cm}|p{12.2cm}|}
\hline
\textbf{CWEs} & CWE-250, CWE-426 \\ \hline
\textbf{CVSS 3.0} & 8.8 \\ \hline
\textbf{Risk Factor} & HIGH \\ \hline
\textbf{Description} & The custom binary \texttt{easysysinfo} operates with SUID root permissions and executes the \texttt{cat} utility via the \texttt{system()} function without specifying an absolute path (e.g., \texttt{/bin/cat}). A local attacker can modify their \texttt{PATH} environment variable to execute a malicious binary named \texttt{cat}, resulting in arbitrary code execution with root privileges. \\ \hline
\textbf{Solution} & \textbf{Fix Untrusted Search Path:} Update the source code of custom executables to specify absolute system paths for all external binaries, or utilize direct system calls (such as \texttt{execve}) instead of spawning a shell with \texttt{system()}. \newline \textbf{Least Privilege:} Audit custom binaries and remove the SUID bit from files that do not strictly require root execution privileges. \\ \hline
\end{tabular}

%%%%%%%%%%%%%%%%%%%%%%%%%%%%%%%%%%%%%%%%%%%%
%
% Apache 2.4.x < 2.4.56 Multiple Vulnerabilities
%
%%%%%%%%%%%%%%%%%%%%%%%%%%%%%%%%%%%%%%%%%%%%
\subsection{Apache 2.4.x < 2.4.56 Multiple Vulnerabilities}

\subsubsection{F-007: HTTP Request Splitting / Smuggling with mod\_rewrite and mod\_proxy}
\begin{tabular}{|p{2.8cm}|p{12.2cm}|}
\hline
\textbf{CVE} & CVE-2023-25690 \\ \hline
\textbf{CVSS 3.0} & 9.8 \\ \hline
\textbf{Risk Factor} & CRITICAL \\ \hline
\textbf{Description} & HTTP request splitting with mod\_rewrite and mod\_proxy: Some mod\_proxy configurations on Apache HTTP Server versions 2.4.0 through 2.4.55 allow a HTTP Request Smuggling attack. Configurations are affected when mod\_proxy is enabled along with some form of RewriteRule or ProxyPassMatch in which a non-specific pattern matches some portion of the user-supplied request-target (URL) data and is then re-inserted into the proxied request-target using variable substitution. Request splitting/smuggling could result in bypass of access controls in the proxy server, proxying unintended URLs to existing origin servers, and cache poisoning. \\ \hline
\textbf{Solution} & Upgrade to Apache version 2.4.56 or later. \\ \hline
\end{tabular}

\subsubsection{F-008: mod\_proxy\_uwsgi HTTP Response Splitting}
\begin{tabular}{|p{2.8cm}|p{12.2cm}|}
\hline
\textbf{CVE} & CVE-2023-27522 \\ \hline
\textbf{CVSS 3.0} & 7.5 \\ \hline
\textbf{Risk Factor} & HIGH \\ \hline
\textbf{Description} & Apache HTTP Server: mod\_proxy\_uwsgi HTTP response splitting: HTTP Response Smuggling vulnerability in Apache HTTP Server via mod\_proxy\_uwsgi. This issue affects Apache HTTP Server: from 2.4.30 through 2.4.55. Special characters in the origin response header can truncate/split the response forwarded to the client. \\ \hline
\textbf{Solution} & Upgrade to Apache version 2.4.56 or later. \\ \hline
\end{tabular}

%%%%%%%%%%%%%%%%%%%%%%%%%%%%%%%%%%%%%%%%%%%%
%
% Apache 2.4.x < 2.4.60 Multiple Vulnerabilities
%
%%%%%%%%%%%%%%%%%%%%%%%%%%%%%%%%%%%%%%%%%%%%
\subsection{Apache 2.4.x < 2.4.60 Multiple Vulnerabilities}

\subsubsection{F-009: WebSocket Protocol Upgrade Null Pointer Dereference}
\begin{tabular}{|p{2.8cm}|p{12.2cm}|}
\hline
\textbf{CVE} & CVE-2024-36387 \\ \hline
\textbf{CVSS 3.0} & 5.4 \\ \hline
\textbf{Risk Factor} & MEDIUM \\ \hline
\textbf{Description} & The version of Apache HTTP Server installed on the remote host is prior to 2.4.60. Serving WebSocket protocol upgrades over a HTTP/2 connection could result in a Null Pointer dereference, leading to a crash of the server process and degraded performance. \\ \hline
\textbf{Solution} & Upgrade to Apache version 2.4.60 or later. \\ \hline
\end{tabular}

\subsubsection{F-010: SSRF in Apache HTTP Server on Windows NTLM Hash Leakage}
\begin{tabular}{|p{2.8cm}|p{12.2cm}|}
\hline
\textbf{CVE} & CVE-2024-38472 \\ \hline
\textbf{CVSS 3.0} & 7.5 \\ \hline
\textbf{Risk Factor} & HIGH \\ \hline
\textbf{Description} & The version of Apache HTTP Server installed on the remote host is prior to 2.4.60. SSRF in Apache HTTP Server on Windows allows potentially leaking NTLM hashes to a malicious server via SSRF and malicious requests or content. Existing configurations accessing UNC paths must configure the UNCList directive to allow access during request processing. \\ \hline
\textbf{Solution} & Upgrade to Apache version 2.4.60 or later. \\ \hline
\end{tabular}

\subsubsection{F-011: mod\_proxy Request URL Encoding Authentication Bypass}
\begin{tabular}{|p{2.8cm}|p{12.2cm}|}
\hline
\textbf{CVE} & CVE-2024-38473 \\ \hline
\textbf{CVSS 3.0} & 8.1 \\ \hline
\textbf{Risk Factor} & HIGH \\ \hline
\textbf{Description} & The version of Apache HTTP Server installed on the remote host is prior to 2.4.60. An encoding issue in mod\_proxy allows request URLs with incorrect encoding to be sent to backend services, potentially bypassing authentication through crafted requests. \\ \hline
\textbf{Solution} & Upgrade to Apache version 2.4.60 or later. \\ \hline
\end{tabular}

\subsubsection{F-012: mod\_rewrite Substitution Encoding Script Execution / Source Disclosure}
\begin{tabular}{|p{2.8cm}|p{12.2cm}|}
\hline
\textbf{CVE} & CVE-2024-38474 \\ \hline
\textbf{CVSS 3.0} & 9.8 \\ \hline
\textbf{Risk Factor} & CRITICAL \\ \hline
\textbf{Description} & The version of Apache HTTP Server installed on the remote host is prior to 2.4.60. A substitution encoding issue in mod\_rewrite allows attackers to execute scripts in directories permitted by configuration but not directly reachable by any URL, or disclose source code of scripts intended only to be executed as CGI. Unsafe RewriteRules may require the UnsafeAllow3F flag. \\ \hline
\textbf{Solution} & Upgrade to Apache version 2.4.60 or later. \\ \hline
\end{tabular}

\subsubsection{F-013: mod\_rewrite Output Escaping File Mapping Vulnerability}
\begin{tabular}{|p{2.8cm}|p{12.2cm}|}
\hline
\textbf{CVE} & CVE-2024-38475 \\ \hline
\textbf{CVSS 3.0} & 9.1 \\ \hline
\textbf{Risk Factor} & CRITICAL \\ \hline
\textbf{Description} & The version of Apache HTTP Server installed on the remote host is prior to 2.4.60. Improper escaping of output in mod\_rewrite allows attackers to map URLs to filesystem locations that are permitted to be served but are not intentionally reachable, resulting in code execution or source code disclosure. Affected substitutions include those using backreferences or variables as the first segment of the substitution. The UnsafePrefixStat flag can be used to opt back in after ensuring substitutions are properly constrained. \\ \hline
\textbf{Solution} & Upgrade to Apache version 2.4.60 or later. \\ \hline
\end{tabular}

\subsubsection{F-014: Core Information Disclosure / SSRF / Script Execution}
\begin{tabular}{|p{2.8cm}|p{12.2cm}|}
\hline
\textbf{CVE} & CVE-2024-38476 \\ \hline
\textbf{CVSS 3.0} & 9.8 \\ \hline
\textbf{Risk Factor} & CRITICAL \\ \hline
\textbf{Description} & The version of Apache HTTP Server installed on the remote host is prior to 2.4.60. A vulnerability in the core of Apache HTTP Server allows information disclosure, SSRF, or local script execution through backend applications with malicious or exploitable response headers. \\ \hline
\textbf{Solution} & Upgrade to Apache version 2.4.60 or later. \\ \hline
\end{tabular}

\subsubsection{F-015: mod\_proxy Null Pointer Dereference Denial of Service}
\begin{tabular}{|p{2.8cm}|p{12.2cm}|}
\hline
\textbf{CVE} & CVE-2024-38477 \\ \hline
\textbf{CVSS 3.0} & 7.5 \\ \hline
\textbf{Risk Factor} & HIGH \\ \hline
\textbf{Description} & The version of Apache HTTP Server installed on the remote host is prior to 2.4.60. A null pointer dereference in mod\_proxy allows an attacker to crash the server through a malicious request. \\ \hline
\textbf{Solution} & Upgrade to Apache version 2.4.60 or later. \\ \hline
\end{tabular}

\subsubsection{F-016: mod\_rewrite Potential SSRF via mod\_proxy}
\begin{tabular}{|p{2.8cm}|p{12.2cm}|}
\hline
\textbf{CVE} & CVE-2024-39573 \\ \hline
\textbf{CVSS 3.0} & 7.5 \\ \hline
\textbf{Risk Factor} & HIGH \\ \hline
\textbf{Description} & The version of Apache HTTP Server installed on the remote host is prior to 2.4.60. A potential SSRF vulnerability in mod\_rewrite allows attackers to cause unsafe RewriteRules to unexpectedly configure URLs to be handled by mod\_proxy. \\ \hline
\textbf{Solution} & Upgrade to Apache version 2.4.60 or later. \\ \hline
\end{tabular}

%%%%%%%%%%%%%%%%%%%%%%%%%%%%%%%%%%%%%%%%%%%%
%
% Apache 2.4.x < 2.4.67 Multiple Vulnerabilities
%
%%%%%%%%%%%%%%%%%%%%%%%%%%%%%%%%%%%%%%%%%%%%
\subsection{Apache 2.4.x < 2.4.67 Multiple Vulnerabilities}

\subsubsection{F-017: HTTP/2 Protocol Double Free and Remote Code Execution}
\begin{tabular}{|p{2.8cm}|p{12.2cm}|}
\hline
\textbf{CVE} & CVE-2026-23918 \\ \hline
\textbf{CVSS 3.0} & 8.8 \\ \hline
\textbf{Risk Factor} & HIGH \\ \hline
\textbf{Description} & The version of Apache httpd installed on the remote host is prior to 2.4.67. Double Free and possible Remote Code Execution vulnerability in Apache HTTP Server with the HTTP/2 protocol. This issue affects Apache HTTP Server version 2.4.66. \\ \hline
\textbf{Solution} & Upgrade to Apache version 2.4.67 or later. \\ \hline
\end{tabular}

\subsubsection{F-018: Local .htaccess Authors Escalation of Privilege}
\begin{tabular}{|p{2.8cm}|p{12.2cm}|}
\hline
\textbf{CVE} & CVE-2026-24072 \\ \hline
\textbf{CVSS 3.0} & 8.8 \\ \hline
\textbf{Risk Factor} & HIGH \\ \hline
\textbf{Description} & The version of Apache httpd installed on the remote host is prior to 2.4.67. An escalation of privilege vulnerability in various modules in Apache HTTP Server 2.4.66 and earlier allows local .htaccess authors to read files with the privileges of the httpd user. \\ \hline
\textbf{Solution} & Upgrade to Apache version 2.4.67 or later. \\ \hline
\end{tabular}

\subsubsection{F-019: mod\_proxy\_ajp Heap-based Buffer Overflow}
\begin{tabular}{|p{2.8cm}|p{12.2cm}|}
\hline
\textbf{CVE} & CVE-2026-28780 \\ \hline
\textbf{CVSS 3.0} & 9.8 \\ \hline
\textbf{Risk Factor} & CRITICAL \\ \hline
\textbf{Description} & The version of Apache httpd installed on the remote host is prior to 2.4.67. Heap-based Buffer Overflow vulnerability in mod\_proxy\_ajp of Apache HTTP Server. If mod\_proxy\_ajp connects to a malicious AJP server, the server can send a malicious AJP message back to mod\_proxy\_ajp and cause it to write 4 attacker-controlled bytes after the end of a heap-based buffer. This issue affects Apache HTTP Server through version 2.4.66. \\ \hline
\textbf{Solution} & Upgrade to Apache version 2.4.67 or later. \\ \hline
\end{tabular}

\subsubsection{F-020: mod\_md Resource Allocation Without Limits via OCSP}
\begin{tabular}{|p{2.8cm}|p{12.2cm}|}
\hline
\textbf{CVE} & CVE-2026-29168 \\ \hline
\textbf{CVSS 3.0} & 7.3 \\ \hline
\textbf{Risk Factor} & HIGH \\ \hline
\textbf{Description} & The version of Apache httpd installed on the remote host is prior to 2.4.67. Allocation of Resources Without Limits or Throttling vulnerability in Apache HTTP Server's mod\_md via OCSP response data. This issue affects Apache HTTP Server from version 2.4.30 through 2.4.66. \\ \hline
\textbf{Solution} & Upgrade to Apache version 2.4.67 or later. \\ \hline
\end{tabular}

\subsubsection{F-021: mod\_dav\_lock NULL Pointer Dereference}
\begin{tabular}{|p{2.8cm}|p{12.2cm}|}
\hline
\textbf{CVE} & CVE-2026-29169 \\ \hline
\textbf{CVSS 3.0} & 7.5 \\ \hline
\textbf{Risk Factor} & HIGH \\ \hline
\textbf{Description} & The version of Apache httpd installed on the remote host is prior to 2.4.67. A NULL pointer dereference vulnerability in mod\_dav\_lock in Apache HTTP Server 2.4.66 and earlier may allow an attacker to crash the server with a malicious request. mod\_dav\_lock is not used internally by mod\_dav or mod\_dav\_fs. The only known use-case for mod\_dav\_lock was mod\_dav\_svn from Apache Subversion earlier than version 1.2.0. \\ \hline
\textbf{Solution} & Upgrade to Apache version 2.4.67 or later. \\ \hline
\end{tabular}

\subsubsection{F-022: Apache 2.4.x < 2.4.67 Vulnerability (CVE-2026-33006)}
\begin{tabular}{|p{2.8cm}|p{12.2cm}|}
\hline
\textbf{CVE} & CVE-2026-33006 \\ \hline
\textbf{CVSS 3.0} & 4.8 \\ \hline
\textbf{Risk Factor} & MEDIUM \\ \hline
\textbf{Description} & A timing attack against mod\_auth\_digest in Apache HTTP Server 2.4.66 allows a bypass of Digest authentication by a remote attacker. \\ \hline
\textbf{Solution} & Upgrade to Apache version 2.4.67 or later. \\ \hline
\end{tabular}

\subsubsection{F-023: Apache 2.4.x < 2.4.67 Vulnerability (CVE-2026-33007)}
\begin{tabular}{|p{2.8cm}|p{12.2cm}|}
\hline
\textbf{CVE} & CVE-2026-33007 \\ \hline
\textbf{CVSS 3.0} & 5.3 \\ \hline
\textbf{Risk Factor} & MEDIUM \\ \hline
\textbf{Description} & A NULL pointer dereference in the mod\_authn\_socache in Apache HTTP Server 2.4.66 and earlier allows an unauthenticated remote user to crash a child process in a caching forward proxy configuration. \\ \hline
\textbf{Solution} & Upgrade to Apache version 2.4.67 or later. \\ \hline
\end{tabular}

\subsubsection{F-024: Apache 2.4.x < 2.4.67 Vulnerability (CVE-2026-33523)}
\begin{tabular}{|p{2.8cm}|p{12.2cm}|}
\hline
\textbf{CVE} & CVE-2026-33523 \\ \hline
\textbf{CVSS 3.0} & 6.5 \\ \hline
\textbf{Risk Factor} & MEDIUM \\ \hline
\textbf{Description} & HTTP response splitting vulnerability in multiple Apache HTTP Server modules with untrusted or compromised backend servers. \\ \hline
\textbf{Solution} & Upgrade to Apache version 2.4.67 or later. \\ \hline
\end{tabular}

\subsubsection{F-025: Apache 2.4.x < 2.4.67 Vulnerability (CVE-2026-33857)}
\begin{tabular}{|p{2.8cm}|p{12.2cm}|}
\hline
\textbf{CVE} & CVE-2026-33857 \\ \hline
\textbf{CVSS 3.0} & 5.3 \\ \hline
\textbf{Risk Factor} & MEDIUM \\ \hline
\textbf{Description} & Out-of-bounds Read vulnerability in mod\_proxy\_ajp of Apache HTTP Server. \\ \hline
\textbf{Solution} & Upgrade to Apache version 2.4.67 or later. \\ \hline
\end{tabular}

\subsubsection{F-026: Apache 2.4.x < 2.4.67 Vulnerability (CVE-2026-34032)}
\begin{tabular}{|p{2.8cm}|p{12.2cm}|}
\hline
\textbf{CVE} & CVE-2026-34032 \\ \hline
\textbf{CVSS 3.0} & 5.3 \\ \hline
\textbf{Risk Factor} & MEDIUM \\ \hline
\textbf{Description} & Improper Null Termination, Out-of-bounds Read vulnerability in Apache HTTP Server. \\ \hline
\textbf{Solution} & Upgrade to Apache version 2.4.67 or later. \\ \hline
\end{tabular}

\subsubsection{F-027: Apache 2.4.x < 2.4.67 Vulnerability (CVE-2026-34059)}
\begin{tabular}{|p{2.8cm}|p{12.2cm}|}
\hline
\textbf{CVE} & CVE-2026-34059 \\ \hline
\textbf{CVSS 3.0} & 7.5 \\ \hline
\textbf{Risk Factor} & HIGH \\ \hline
\textbf{Description} & Buffer Over-read vulnerability in Apache HTTP Server. \\ \hline
\textbf{Solution} & Upgrade to Apache version 2.4.67 or later. \\ \hline
\end{tabular}

%%%%%%%%%%%%%%%%%%%%%%%%%%%%%%%%%%%%%%%%%%%%
%
% Apache 2.4.x < 2.4.47 Multiple Vulnerabilities
%
%%%%%%%%%%%%%%%%%%%%%%%%%%%%%%%%%%%%%%%%%%%%
\subsection{Apache 2.4.x < 2.4.47 Multiple Vulnerabilities}

\subsubsection{F-028: mod\_proxy\_wstunnel and mod\_proxy\_http Upgradable Protocols Negotiation}
\begin{tabular}{|p{2.8cm}|p{12.2cm}|}
\hline
\textbf{CVE} & CVE-2019-17567 \\ \hline
\textbf{CVSS 3.0} & 5.3 \\ \hline
\textbf{Risk Factor} & MEDIUM \\ \hline
\textbf{Description} & Apache HTTP Server versions 2.4.6 to 2.4.46 mod\_proxy\_wstunnel configured on an URL that is not necessarily Upgraded by the origin server was tunneling the whole connection regardless, thus allowing for subsequent requests on the same connection to pass through with no HTTP validation, authentication or authorization possibly configured. \\ \hline
\textbf{Solution} & Upgrade to Apache version 2.4.47 or later. \\ \hline
\end{tabular}

\subsubsection{F-029: Windows Local Process Stopping Vulnerability}
\begin{tabular}{|p{2.8cm}|p{12.2cm}|}
\hline
\textbf{CVE} & CVE-2020-13938 \\ \hline
\textbf{CVSS 3.0} & 5.5 \\ \hline
\textbf{Risk Factor} & MEDIUM \\ \hline
\textbf{Description} & Unprivileged local users can stop httpd on Windows. \\ \hline
\textbf{Solution} & Upgrade to Apache version 2.4.47 or later. \\ \hline
\end{tabular}

\subsubsection{F-030: mod\_proxy\_http NULL Pointer Dereference Denial of Service}
\begin{tabular}{|p{2.8cm}|p{12.2cm}|}
\hline
\textbf{CVE} & CVE-2020-13950 \\ \hline
\textbf{CVSS 3.0} & 7.5 \\ \hline
\textbf{Risk Factor} & HIGH \\ \hline
\textbf{Description} & Apache HTTP Server versions 2.4.41 to 2.4.46 mod\_proxy\_http can be made to crash (NULL pointer dereference) with specially crafted requests using both Content-Length and Transfer-Encoding headers, leading to a Denial of Service. \\ \hline
\textbf{Solution} & Upgrade to Apache version 2.4.47 or later. \\ \hline
\end{tabular}

\subsubsection{F-031: mod\_auth\_digest Stack Overflow via Digest Nonce}
\begin{tabular}{|p{2.8cm}|p{12.2cm}|}
\hline
\textbf{CVE} & CVE-2020-35452 \\ \hline
\textbf{CVSS 3.0} & 7.3 \\ \hline
\textbf{Risk Factor} & HIGH \\ \hline
\textbf{Description} & A specially crafted Digest nonce can cause a stack overflow in mod\_auth\_digest. There is no report of this overflow being exploitable, nor the Apache HTTP Server team could create one, though some particular compiler and/or compilation option might make it possible, with limited consequences anyway due to the size (a single byte) and the value (zero byte) of the overflow. \\ \hline
\textbf{Solution} & Upgrade to Apache version 2.4.47 or later. \\ \hline
\end{tabular}

\subsubsection{F-032: mod\_session NULL Pointer Dereference Denial of Service}
\begin{tabular}{|p{2.8cm}|p{12.2cm}|}
\hline
\textbf{CVE} & CVE-2021-26690 \\ \hline
\textbf{CVSS 3.0} & 7.5 \\ \hline
\textbf{Risk Factor} & HIGH \\ \hline
\textbf{Description} & A specially crafted Cookie header handled by mod\_session can cause a NULL pointer dereference and crash, leading to a possible Denial Of Service. \\ \hline
\textbf{Solution} & Upgrade to Apache version 2.4.47 or later. \\ \hline
\end{tabular}

\subsubsection{F-033: mod\_session NULL Pointer Dereference DoS via SessionHeader}
\begin{tabular}{|p{2.8cm}|p{12.2cm}|}
\hline
\textbf{CVE} & CVE-2021-26691 \\ \hline
\textbf{CVSS 3.0} & 9.8 \\ \hline
\textbf{Risk Factor} & CRITICAL \\ \hline
\textbf{Description} & In Apache HTTP Server versions 2.4.0 to 2.4.46 a specially crafted SessionHeader sent by an origin server could cause a heap overflow. \\ \hline
\textbf{Solution} & Upgrade to Apache version 2.4.47 or later. \\ \hline
\end{tabular}

\subsubsection{F-034: Unexpected <Location> Section Matching with 'MergeSlashes OFF'}
\begin{tabular}{|p{2.8cm}|p{12.2cm}|}
\hline
\textbf{CVE} & CVE-2021-30641 \\ \hline
\textbf{CVSS 3.0} & 5.3 \\ \hline
\textbf{Risk Factor} & MEDIUM \\ \hline
\textbf{Description} & The version of Apache httpd installed on the remote host is prior to 2.4.47. Unexpected <Location> section matching with 'MergeSlashes OFF'. \\ \hline
\textbf{Solution} & Upgrade to Apache version 2.4.47 or later. \\ \hline
\end{tabular}

%%%%%%%%%%%%%%%%%%%%%%%%%%%%%%%%%%%%%%%%%%%%
%
% Apache 2.4.x < 2.4.52 mod_lua Buffer Overflow
%
%%%%%%%%%%%%%%%%%%%%%%%%%%%%%%%%%%%%%%%%%%%%
\subsection{Apache 2.4.x < 2.4.52 mod\_lua Buffer Overflow}

\subsubsection{F-035: Apache 2.4.x < 2.4.52 mod\_lua Buffer Overflow}
\begin{tabular}{|p{2.8cm}|p{12.2cm}|}
\hline
\textbf{CVE} & CVE-2021-44790 \\ \hline
\textbf{CVSS 3.0} & 9.8 \\ \hline
\textbf{Risk Factor} & CRITICAL \\ \hline
\textbf{Description} & The version of Apache httpd installed on the remote host is prior to 2.4.52. It is, therefore, affected by a flaw related to mod\_lua when handling multipart content. A carefully crafted request body can cause a buffer overflow in the mod\_lua multipart parser (r:parsebody() called from Lua scripts). The Apache httpd team is not aware of an exploit for the vulnerability though it might be possible to craft one. \\ \hline
\textbf{Solution} & Upgrade to Apache version 2.4.52 or later. \\ \hline
\end{tabular}

%%%%%%%%%%%%%%%%%%%%%%%%%%%%%%%%%%%%%%%%%%%%
%
% Apache 2.4.x < 2.4.53 Multiple Vulnerabilities
%
%%%%%%%%%%%%%%%%%%%%%%%%%%%%%%%%%%%%%%%%%%%%
\subsection{Apache 2.4.x < 2.4.53 Multiple Vulnerabilities}

\subsubsection{F-036: mod\_lua Use of Uninitialized Value in r:parsebody}
\begin{tabular}{|p{2.8cm}|p{12.2cm}|}
\hline
\textbf{CVE} & CVE-2022-22719 \\ \hline
\textbf{CVSS 3.0} & 7.5 \\ \hline
\textbf{Risk Factor} & HIGH \\ \hline
\textbf{Description} & The version of Apache HTTP Server installed on the remote host is prior to 2.4.53. mod\_lua Use of uninitialized value of in r:parsebody: A carefully crafted request body can cause a read to a random memory area which could cause the process to crash. This issue affects Apache HTTP Server 2.4.52 and earlier. \\ \hline
\textbf{Solution} & Upgrade to Apache version 2.4.53 or later. \\ \hline
\end{tabular}

\subsubsection{F-037: HTTP Request Smuggling Connection Closing Issue}
\begin{tabular}{|p{2.8cm}|p{12.2cm}|}
\hline
\textbf{CVE} & CVE-2022-22720 \\ \hline
\textbf{CVSS 3.0} & 9.8 \\ \hline
\textbf{Risk Factor} & CRITICAL \\ \hline
\textbf{Description} & The version of Apache HTTP Server installed on the remote host is prior to 2.4.53. HTTP request smuggling: Apache HTTP Server 2.4.52 and earlier fails to close inbound connection when errors are encountered while discarding the request body, exposing the server to HTTP Request Smuggling attacks. \\ \hline
\textbf{Solution} & Upgrade to Apache version 2.4.53 or later. \\ \hline
\end{tabular}

\subsubsection{F-038: LimitXMLRequestBody Integer Overflow / Out-of-bounds Write}
\begin{tabular}{|p{2.8cm}|p{12.2cm}|}
\hline
\textbf{CVE} & CVE-2022-22721 \\ \hline
\textbf{CVSS 3.0} & 9.1 \\ \hline
\textbf{Risk Factor} & CRITICAL \\ \hline
\textbf{Description} & The version of Apache HTTP Server installed on the remote host is prior to 2.4.53. Possible buffer overflow with very large or unlimited LimitXMLRequestBody in core: If LimitXMLRequestBody is set to allow request bodies larger than 350MB (defaults to 1M) on 32-bit systems, an integer overflow can occur, later causing out-of-bounds writes. This issue affects Apache HTTP Server 2.4.52 and earlier. \\ \hline
\textbf{Solution} & Upgrade to Apache version 2.4.53 or later. \\ \hline
\end{tabular}

\subsubsection{F-039: mod\_sed Read/Write Beyond Bounds Heap Overwrite}
\begin{tabular}{|p{2.8cm}|p{12.2cm}|}
\hline
\textbf{CVE} & CVE-2022-23943 \\ \hline
\textbf{CVSS 3.0} & 9.8 \\ \hline
\textbf{Risk Factor} & CRITICAL \\ \hline
\textbf{Description} & The version of Apache HTTP Server installed on the remote host is prior to 2.4.53. Read/write beyond bounds in mod\_sed: An out-of-bounds write vulnerability in mod\_sed of Apache HTTP Server allows an attacker to overwrite heap memory with potentially attacker-provided data. This issue affects Apache HTTP Server 2.4.52 and prior versions. \\ \hline
\textbf{Solution} & Upgrade to Apache version 2.4.53 or later. \\ \hline
\end{tabular}

%%%%%%%%%%%%%%%%%%%%%%%%%%%%%%%%%%%%%%%%%%%%
%
% Apache 2.4.x < 2.4.54 Authentication Bypass
%
%%%%%%%%%%%%%%%%%%%%%%%%%%%%%%%%%%%%%%%%%%%%
\subsection{Apache 2.4.x < 2.4.54 Authentication Bypass}

\subsubsection{F-040: Apache 2.4.x < 2.4.54 Authentication Bypass via mod\_proxy}
\begin{tabular}{|p{2.8cm}|p{12.2cm}|}
\hline
\textbf{CVE} & CVE-2022-31813 \\ \hline
\textbf{CVSS 3.0} & 9.8 \\ \hline
\textbf{Risk Factor} & CRITICAL \\ \hline
\textbf{Description} & The version of Apache httpd installed on the remote host is prior to 2.4.54. It is, therefore, affected by an authentication bypass vulnerability as referenced in the 2.4.54 advisory. X-Forwarded-For dropped by hop-by-hop mechanism in mod\_proxy: Apache HTTP Server 2.4.53 and earlier may not send the X-Forwarded-* headers to the origin server based on client side Connection header hop-by-hop mechanism. This may be used to bypass IP based authentication on the origin server/application. \\ \hline
\textbf{Solution} & Upgrade to Apache version 2.4.54 or later. \\ \hline
\end{tabular}

%%%%%%%%%%%%%%%%%%%%%%%%%%%%%%%%%%%%%%%%%%%%
%
% Apache 2.4.x < 2.4.55 Multiple Vulnerabilities
%
%%%%%%%%%%%%%%%%%%%%%%%%%%%%%%%%%%%%%%%%%%%%
\subsection{Apache 2.4.x < 2.4.55 Multiple Vulnerabilities}

\subsubsection{F-041: If: Request Header Heap Memory Read/Write}
\begin{tabular}{|p{2.8cm}|p{12.2cm}|}
\hline
\textbf{CVE} & CVE-2006-20001 \\ \hline
\textbf{CVSS 3.0} & 7.5 \\ \hline
\textbf{Risk Factor} & HIGH \\ \hline
\textbf{Description} & The version of Apache httpd installed on the remote host is prior to 2.4.55. A carefully crafted If: request header can cause a memory read, or write of a single zero byte, in a pool (heap) memory location beyond the header value sent. This could cause the process to crash. This issue affects Apache HTTP Server 2.4.54 and earlier. \\ \hline
\textbf{Solution} & Upgrade to Apache version 2.4.55 or later. \\ \hline
\end{tabular}

\subsubsection{F-042: mod\_proxy\_ajp HTTP Request Smuggling}
\begin{tabular}{|p{2.8cm}|p{12.2cm}|}
\hline
\textbf{CVE} & CVE-2022-36760 \\ \hline
\textbf{CVSS 3.0} & 9.0 \\ \hline
\textbf{Risk Factor} & CRITICAL \\ \hline
\textbf{Description} & The version of Apache httpd installed on the remote host is prior to 2.4.55. Inconsistent Interpretation of HTTP Requests ('HTTP Request Smuggling') vulnerability in mod\_proxy\_ajp of Apache HTTP Server allows an attacker to smuggle requests to the AJP server it forwards requests to. This issue affects Apache HTTP Server 2.4.54 and prior versions. \\ \hline
\textbf{Solution} & Upgrade to Apache version 2.4.55 or later. \\ \hline
\end{tabular}

\subsubsection{F-043: Backend Response Header Early Truncation}
\begin{tabular}{|p{2.8cm}|p{12.2cm}|}
\hline
\textbf{CVE} & CVE-2022-37436 \\ \hline
\textbf{CVSS 3.0} & 5.3 \\ \hline
\textbf{Risk Factor} & MEDIUM \\ \hline
\textbf{Description} & The version of Apache httpd installed on the remote host is prior to 2.4.55. Prior to Apache HTTP Server 2.4.55, a malicious backend can cause the response headers to be truncated early, resulting in some headers being incorporated into the response body. If the later headers have any security purpose, they will not be interpreted by the client. \\ \hline
\textbf{Solution} & Upgrade to Apache version 2.4.55 or later. \\ \hline
\end{tabular}

%%%%%%%%%%%%%%%%%%%%%%%%%%%%%%%%%%%%%%%%%%%%
%
% Apache 2.4.x < 2.4.58 Multiple Vulnerabilities
%
%%%%%%%%%%%%%%%%%%%%%%%%%%%%%%%%%%%%%%%%%%%%
\subsection{Apache 2.4.x < 2.4.58 Multiple Vulnerabilities}

\subsubsection{F-044: HTTP/2 Connection Block DoS via Initial Window Size 0}
\begin{tabular}{|p{2.8cm}|p{12.2cm}|}
\hline
\textbf{CVE} & CVE-2023-43622 \\ \hline
\textbf{CVSS 3.0} & 7.5 \\ \hline
\textbf{Risk Factor} & HIGH \\ \hline
\textbf{Description} & The version of Apache httpd installed on the remote host is prior to 2.4.58. Apache HTTP Server: DoS in HTTP/2 with initial windows size 0: An attacker opening an HTTP/2 connection with an initial window size of 0 was able to block handling of that connection indefinitely in Apache HTTP Server. This could be used to exhaust worker resources in the server, similar to the well-known slow loris attack pattern. This issue affects Apache HTTP Server versions from 2.4.55 through 2.4.57. \\ \hline
\textbf{Solution} & Upgrade to Apache version 2.4.58 or later. \\ \hline
\end{tabular}

\subsubsection{F-045: HTTP/2 Stream Deferred Memory Release DoS}
\begin{tabular}{|p{2.8cm}|p{12.2cm}|}
\hline
\textbf{CVE} & CVE-2023-45802 \\ \hline
\textbf{CVSS 3.0} & 5.9 \\ \hline
\textbf{Risk Factor} & MEDIUM \\ \hline
\textbf{Description} & The version of Apache httpd installed on the remote host is prior to 2.4.58. Apache HTTP Server: HTTP/2 stream memory not reclaimed right away on RST: When an HTTP/2 stream was reset (RST frame) by a client, there was a time window where the request's memory resources were not reclaimed immediately. Instead, de-allocation was deferred to connection close. A client could send new requests and resets, keeping the connection busy and open and causing the memory footprint to continue growing. The process might run out of memory before connection closure. This issue was identified during testing of CVE-2023-44487 (HTTP/2 Rapid Reset Exploit). During normal HTTP/2 use, the probability of triggering this issue is very low. \\ \hline
\textbf{Solution} & Upgrade to Apache version 2.4.58 or later. \\ \hline
\end{tabular}

%%%%%%%%%%%%%%%%%%%%%%%%%%%%%%%%%%%%%%%%%%%%
%
% Apache 2.4.x < 2.4.58 Out-of-Bounds Read (CVE-2023-31122)
%
%%%%%%%%%%%%%%%%%%%%%%%%%%%%%%%%%%%%%%%%%%%%
\subsection{Apache 2.4.x < 2.4.58 Out-of-Bounds Read (CVE-2023-31122)}

\subsubsection{F-046: Apache 2.4.x < 2.4.58 Out-of-Bounds Read (CVE-2023-31122)}
\begin{tabular}{|p{2.8cm}|p{12.2cm}|}
\hline
\textbf{CVE} & CVE-2023-31122 \\ \hline
\textbf{CVSS 3.0} & 7.5 \\ \hline
\textbf{Risk Factor} & HIGH \\ \hline
\textbf{Description} & The version of Apache httpd installed on the remote host is prior to 2.4.58. It is, therefore, affected by an out-of-bounds read vulnerability as referenced in the 2.4.58 advisory. Apache HTTP Server mod\_macro buffer over-read: An out-of-bounds Read vulnerability exists in mod\_macro of Apache HTTP Server. This issue affects Apache HTTP Server through version 2.4.57. \\ \hline
\textbf{Solution} & Upgrade to Apache version 2.4.58 or later. \\ \hline
\end{tabular}

%%%%%%%%%%%%%%%%%%%%%%%%%%%%%%%%%%%%%%%%%%%%
%
% Apache 2.4.x < 2.4.59 Multiple Vulnerabilities
%
%%%%%%%%%%%%%%%%%%%%%%%%%%%%%%%%%%%%%%%%%%%%
\subsection{Apache 2.4.x < 2.4.59 Multiple Vulnerabilities}

\subsubsection{F-047: HTTP Response Splitting in Multiple Modules}
\begin{tabular}{|p{2.8cm}|p{12.2cm}|}
\hline
\textbf{CVE} & CVE-2023-38709 \\ \hline
\textbf{CVSS 3.0} & 7.3 \\ \hline
\textbf{Risk Factor} & HIGH \\ \hline
\textbf{Description} & The version of Apache httpd installed on the remote host is prior to 2.4.59. It is, therefore, affected by an HTTP Response Splitting vulnerability in multiple modules as referenced in the 2.4.59 advisory. \\ \hline
\textbf{Solution} & Upgrade to Apache version 2.4.59 or later. \\ \hline
\end{tabular}

\subsubsection{F-048: HTTP Response Splitting Desynchronization Attack}
\begin{tabular}{|p{2.8cm}|p{12.2cm}|}
\hline
\textbf{CVE} & CVE-2024-24795 \\ \hline
\textbf{CVSS 3.0} & 6.3 \\ \hline
\textbf{Risk Factor} & MEDIUM \\ \hline
\textbf{Description} & The version of Apache httpd installed on the remote host is prior to 2.4.59. Apache HTTP Server: HTTP Response Splitting in multiple modules: HTTP response splitting in multiple modules of Apache HTTP Server allows an attacker that can inject malicious response headers into backend applications to cause an HTTP desynchronization attack. This issue is fixed in version 2.4.59. \\ \hline
\textbf{Solution} & Upgrade to Apache version 2.4.59 or later. \\ \hline
\end{tabular}

\subsubsection{F-049: HTTP/2 DoS by Memory Exhaustion on Continuation Frames}
\begin{tabular}{|p{2.8cm}|p{12.2cm}|}
\hline
\textbf{CVE} & CVE-2024-27316 \\ \hline
\textbf{CVSS 3.0} & 7.5 \\ \hline
\textbf{Risk Factor} & HIGH \\ \hline
\textbf{Description} & The version of Apache httpd installed on the remote host is prior to 2.4.59. Apache HTTP Server: HTTP/2 DoS by memory exhaustion on endless continuation frames: HTTP/2 incoming headers exceeding the configured limit are temporarily buffered in nghttp2 in order to generate an informative HTTP 413 response. If a client does not stop sending headers, this can lead to memory exhaustion. \\ \hline
\textbf{Solution} & Upgrade to Apache version 2.4.59 or later. \\ \hline
\end{tabular}

%%%%%%%%%%%%%%%%%%%%%%%%%%%%%%%%%%%%%%%%%%%%
%
% Apache 2.4.x < 2.4.68 Multiple Vulnerabilities
%
%%%%%%%%%%%%%%%%%%%%%%%%%%%%%%%%%%%%%%%%%%%%
\subsection{Apache 2.4.x < 2.4.68 Multiple Vulnerabilities}

\subsubsection{F-050: Apache 2.4.x < 2.4.68 Vulnerability (CVE-2026-29167)}
\begin{tabular}{|p{2.8cm}|p{12.2cm}|}
\hline
\textbf{CVE} & CVE-2026-29167 \\ \hline
\textbf{CVSS 3.0} & 9.8 \\ \hline
\textbf{Risk Factor} & CRITICAL \\ \hline
\textbf{Description} & Use After Free vulnerability in Apache HTTP Server with mod\_ldap in per-directory configuration This issue affects Apache HTTP Server: from 2.4.0 through 2.4.67. \\ \hline
\textbf{Solution} & Upgrade to Apache version 2.4.68 or later. \\ \hline
\end{tabular}

\subsubsection{F-051: Apache 2.4.x < 2.4.68 Vulnerability (CVE-2026-29170)}
\begin{tabular}{|p{2.8cm}|p{12.2cm}|}
\hline
\textbf{CVE} & CVE-2026-29170 \\ \hline
\textbf{CVSS 3.0} & 6.1 \\ \hline
\textbf{Risk Factor} & MEDIUM \\ \hline
\textbf{Description} & A cross-site scripting vulnerability exists in mod\_proxy\_ftp's HTML directory list generation in Apache HTTP Server 2.4.67 and earlier when listing FTP directory contents either via forward or reverse proxy configuration. \\ \hline
\textbf{Solution} & Upgrade to Apache version 2.4.68 or later. \\ \hline
\end{tabular}

\subsubsection{F-052: Apache 2.4.x < 2.4.68 Vulnerability (CVE-2026-34355)}
\begin{tabular}{|p{2.8cm}|p{12.2cm}|}
\hline
\textbf{CVE} & CVE-2026-34355 \\ \hline
\textbf{CVSS 3.0} & 7.5 \\ \hline
\textbf{Risk Factor} & HIGH \\ \hline
\textbf{Description} & A buffer overflow in mod\_proxy\_html in Apache HTTP Server 2.4.67 and earlier allows an attack by an untrusted backend. \\ \hline
\textbf{Solution} & Upgrade to Apache version 2.4.68 or later. \\ \hline
\end{tabular}

\subsubsection{F-053: Apache 2.4.x < 2.4.68 Vulnerability (CVE-2026-34356)}
\begin{tabular}{|p{2.8cm}|p{12.2cm}|}
\hline
\textbf{CVE} & CVE-2026-34356 \\ \hline
\textbf{CVSS 3.0} & 7.5 \\ \hline
\textbf{Risk Factor} & HIGH \\ \hline
\textbf{Description} & Heap-based buffer overflow in Apache HTTP Server is triggered by malicious backend servers using ProxyPassReverseCookie* directives. The flaw can corrupt data on the server’s heap, potentially allowing an attacker to compromise the integrity of the HTTP server process. \\ \hline
\textbf{Solution} & Upgrade to Apache version 2.4.68 or later. \\ \hline
\end{tabular}

\subsubsection{F-054: Apache 2.4.x < 2.4.68 Vulnerability (CVE-2026-42535)}
\begin{tabular}{|p{2.8cm}|p{12.2cm}|}
\hline
\textbf{CVE} & CVE-2026-42535 \\ \hline
\textbf{CVSS 3.0} & 9.1 \\ \hline
\textbf{Risk Factor} & CRITICAL \\ \hline
\textbf{Description} & A path handling flaw in mod\_dav\_fs of Apache HTTP Server up to version 2.4.67 allows an authenticated WebDAV content author to directly manipulate trusted DAV property databases. \\ \hline
\textbf{Solution} & Upgrade to Apache version 2.4.68 or later. \\ \hline
\end{tabular}

\subsubsection{F-055: mod\_xml2enc Heap Buffer Overflow}
\begin{tabular}{|p{2.8cm}|p{12.2cm}|}
\hline
\textbf{CVE} & CVE-2026-42536 \\ \hline
\textbf{CVSS 3.0} & 7.5 \\ \hline
\textbf{Risk Factor} & HIGH \\ \hline
\textbf{Description} & A heap-based buffer overflow occurs when Apache HTTP Server processes untrusted XML content through mod\_xml2enc, specifically during the xml2StartParse function. The vulnerability could allow an attacker to manipulate allocated memory on the heap, potentially leading to application crashes or the execution of arbitrary code. The unsafe handling of untrusted XML data enables the attacker to trigger the overflow during normal HTTP request processing, making the flaw exploitable remotely. \\ \hline
\textbf{Solution} & Upgrade to Apache version 2.4.68 or later. \\ \hline
\end{tabular}

\subsubsection{F-056: Out-of-bounds Read in \texttt{merge\_response\_headers}}
\begin{tabular}{|p{2.8cm}|p{12.2cm}|}
\hline
\textbf{CVE} & CVE-2026-43951 \\ \hline
\textbf{CVSS 3.0} & 6.5 \\ \hline
\textbf{Risk Factor} & MEDIUM \\ \hline
\textbf{Description} & Out-of-bounds Read in \texttt{merge\_response\_headers}: An out-of-bounds read vulnerability exists in Apache HTTP Server with mod\_headers and mod\_mime when handling multiple response languages. This issue affects Apache HTTP Server versions from 2.4.0 through 2.4.67 and may cause a crash. \\ \hline
\textbf{Solution} & Upgrade to Apache version 2.4.68 or later. \\ \hline
\end{tabular}

\subsubsection{F-057: Escalation of Privilege via Expressions in .htaccess}
\begin{tabular}{|p{2.8cm}|p{12.2cm}|}
\hline
\textbf{CVE} & CVE-2026-44119 \\ \hline
\textbf{CVSS 3.0} & 5.5 \\ \hline
\textbf{Risk Factor} & MEDIUM \\ \hline
\textbf{Description} & Escalation of privilege through expressions in .htaccess in multiple modules: An improper privilege management vulnerability in Apache HTTP Server 2.4.67 and earlier allows local .htaccess authors to read files with the privileges of the httpd user. \\ \hline
\textbf{Solution} & Upgrade to Apache version 2.4.68 or later. \\ \hline
\end{tabular}

\subsubsection{F-058: Stack Buffer Over-Read in mod\_ssl OCSP \texttt{send\_request}}
\begin{tabular}{|p{2.8cm}|p{12.2cm}|}
\hline
\textbf{CVE} & CVE-2026-44185 \\ \hline
\textbf{CVSS 3.0} & 7.3 \\ \hline
\textbf{Risk Factor} & HIGH \\ \hline
\textbf{Description} & Stack Buffer Over-Read in mod\_ssl OCSP \texttt{send\_request}: A buffer over-read vulnerability exists in Apache HTTP Server through outbound OCSP requests to an attacker-controlled OCSP server. This issue affects Apache HTTP Server versions from 2.4.0 through 2.4.67. \\ \hline
\textbf{Solution} & Upgrade to Apache version 2.4.68 or later. \\ \hline
\end{tabular}

\subsubsection{F-059: Apache 2.4.x < 2.4.68 Vulnerability (CVE-2026-44186)}
\begin{tabular}{|p{2.8cm}|p{12.2cm}|}
\hline
\textbf{CVE} & CVE-2026-44186 \\ \hline
\textbf{CVSS 3.0} & 7.3 \\ \hline
\textbf{Risk Factor} & HIGH \\ \hline
\textbf{Description} & Loop with Unreachable Exit Condition ('Infinite Loop') vulnerability in the mod\_proxy\_ftp module in Apache HTTP Server with an attacker controlled backend FTP server. \\ \hline
\textbf{Solution} & Upgrade to Apache version 2.4.68 or later. \\ \hline
\end{tabular}

\subsubsection{F-060: Apache 2.4.x < 2.4.68 Vulnerability (CVE-2026-44631)}
\begin{tabular}{|p{2.8cm}|p{12.2cm}|}
\hline
\textbf{CVE} & CVE-2026-44631 \\ \hline
\textbf{CVSS 3.0} & 9.8 \\ \hline
\textbf{Risk Factor} & CRITICAL \\ \hline
\textbf{Description} & The vulnerability is a buffer underflow caused by a signed character overflow in the ap\_regname function, which processes regular expression patterns in Apache HTTP Server configuration files. \\ \hline
\textbf{Solution} & Upgrade to Apache version 2.4.68 or later. \\ \hline
\end{tabular}

\subsubsection{F-061: Apache 2.4.x < 2.4.68 Vulnerability (CVE-2026-48913)}
\begin{tabular}{|p{2.8cm}|p{12.2cm}|}
\hline
\textbf{CVE} & CVE-2026-48913 \\ \hline
\textbf{CVSS 3.0} & 7.3 \\ \hline
\textbf{Risk Factor} & HIGH \\ \hline
\textbf{Description} & Use After Free vulnerability in Apache HTTP Server module mod\_http2 when file handles are already exhausted. \\ \hline
\textbf{Solution} & Upgrade to Apache version 2.4.68 or later. \\ \hline
\end{tabular}

\subsubsection{F-062: HTTP/2 Bomb Remote Denial-of-Service}
\begin{tabular}{|p{2.8cm}|p{12.2cm}|}
\hline
\textbf{CVE} & CVE-2026-49975 \\ \hline
\textbf{CVSS 3.0} & 7.5 \\ \hline
\textbf{Risk Factor} & HIGH \\ \hline
\textbf{Description} & CVE-2026-49975, also known as HTTP/2 Bomb, is a remote denial-of-service vulnerability affecting major web servers, including Apache HTTP Server. The vulnerable behavior exists in the default HTTP/2 configuration and can be exploited to exhaust server resources. \\ \hline
\textbf{Solution} & Upgrade to Apache version 2.4.68 or later. \\ \hline
\end{tabular}

%%%%%%%%%%%%%%%%%%%%%%%%%%%%%%%%%%%%%%%%%%%%
%
% Apache 2.4.x >= 2.4.7 / < 2.4.52 Forward Proxy DoS / SSRF
%
%%%%%%%%%%%%%%%%%%%%%%%%%%%%%%%%%%%%%%%%%%%%
\subsection{Apache 2.4.x >= 2.4.7 / < 2.4.52 Forward Proxy DoS / SSRF}

\subsubsection{F-063: Forward Proxy NULL Pointer Dereference / SSRF}
\begin{tabular}{|p{2.8cm}|p{12.2cm}|}
\hline
\textbf{CVE} & CVE-2021-44224 \\ \hline
\textbf{CVSS 3.0} & 8.2 \\ \hline
\textbf{Risk Factor} & HIGH \\ \hline
\textbf{Description} & A crafted URI sent to Apache HTTP Server configured as a forward proxy (\texttt{ProxyRequests On}) can cause a crash through a NULL pointer dereference. Additionally, configurations mixing forward and reverse proxy declarations may allow requests to be directed to a declared Unix Domain Socket endpoint, resulting in a Server-Side Request Forgery (SSRF) vulnerability. \\ \hline
\textbf{Solution} & Upgrade to Apache version 2.4.52 or later. \\ \hline
\end{tabular}

\subsubsection{F-064: mod\_lua Buffer Overflow in Multipart Content}
\begin{tabular}{|p{2.8cm}|p{12.2cm}|}
\hline
\textbf{CVE} & CVE-2021-44790 \\ \hline
\textbf{CVSS 3.0} & 9.8 \\ \hline
\textbf{Risk Factor} & CRITICAL \\ \hline
\textbf{Description} & A carefully crafted request body can cause a buffer overflow in the mod\_lua multipart parser (r:parsebody() called from Lua scripts). The Apache httpd team is not aware of an exploit for the vulnerabilty though it might be possible to craft one. This issue affects Apache HTTP Server 2.4.51 and earlier. \\ \hline
\textbf{Solution} & Upgrade to Apache version 2.4.52 or later. \\ \hline
\end{tabular}

%%%%%%%%%%%%%%%%%%%%%%%%%%%%%%%%%%%%%%%%%%%%
%
% Apache < 2.4.49 Multiple Vulnerabilities
%
%%%%%%%%%%%%%%%%%%%%%%%%%%%%%%%%%%%%%%%%%%%%
\subsection{Apache < 2.4.49 Multiple Vulnerabilities}

\subsubsection{F-065: Malformed HTTP Request NULL Pointer Dereference DoS}
\begin{tabular}{|p{2.8cm}|p{12.2cm}|}
\hline
\textbf{CVE} & CVE-2021-34798 \\ \hline
\textbf{CVSS 3.0} & 7.5 \\ \hline
\textbf{Risk Factor} & HIGH \\ \hline
\textbf{Description} & Malformed HTTP requests may cause the server to dereference a NULL pointer, resulting in a denial of service condition. \\ \hline
\textbf{Solution} & Upgrade to Apache version 2.4.49 or later. \\ \hline
\end{tabular}

\subsubsection{F-066: ap\_escape\_quotes() Buffer Overwrite}
\begin{tabular}{|p{2.8cm}|p{12.2cm}|}
\hline
\textbf{CVE} & CVE-2021-39275 \\ \hline
\textbf{CVSS 3.0} & 9.8 \\ \hline
\textbf{Risk Factor} & CRITICAL \\ \hline
\textbf{Description} & The \texttt{ap\_escape\_quotes()} function may write beyond the end of a buffer when given malicious input. No included Apache modules pass untrusted data to these functions, but third-party or external modules may be affected. \\ \hline
\textbf{Solution} & Upgrade to Apache version 2.4.49 or later. \\ \hline
\end{tabular}

%%%%%%%%%%%%%%%%%%%%%%%%%%%%%%%%%%%%%%%%%%%%
%
% Apache 2.4.x < 2.4.54 HTTP Request Smuggling Vulnerability
%
%%%%%%%%%%%%%%%%%%%%%%%%%%%%%%%%%%%%%%%%%%%%
\subsection{Apache 2.4.x < 2.4.54 HTTP Request Smuggling Vulnerability}

\subsubsection{F-067: mod\_proxy\_ajp HTTP Request Smuggling}
\begin{tabular}{|p{2.8cm}|p{12.2cm}|}
\hline
\textbf{CVE} & CVE-2022-26377 \\ \hline
\textbf{CVSS 3.0} & 6.5 \\ \hline
\textbf{Risk Factor} & MEDIUM \\ \hline
\textbf{Description} & The version of Apache httpd installed on the remote host is prior to 2.4.54. It is, therefore, affected by an HTTP request smuggling vulnerability as referenced in the 2.4.54 advisory. Possible request smuggling in mod\_proxy\_ajp: An Inconsistent Interpretation of HTTP Requests ('HTTP Request Smuggling') vulnerability exists in mod\_proxy\_ajp of Apache HTTP Server. An attacker may exploit this vulnerability to smuggle requests to the AJP server to which Apache forwards requests. This issue affects Apache HTTP Server version 2.4.53 and prior versions. \\ \hline
\textbf{Solution} & Upgrade to Apache version 2.4.54 or later. \\ \hline
\end{tabular}

%%%%%%%%%%%%%%%%%%%%%%%%%%%%%%%%%%%%%%%%%%%%
%
% Apache 2.4.x < 2.4.54 Multiple Vulnerabilities
%
%%%%%%%%%%%%%%%%%%%%%%%%%%%%%%%%%%%%%%%%%%%%
\subsection{Apache 2.4.x < 2.4.54 Multiple Vulnerabilities}

\subsubsection{F-068: Read Beyond Bounds via \texttt{ap\_rwrite()}}
\begin{tabular}{|p{2.8cm}|p{12.2cm}|}
\hline
\textbf{CVE} & CVE-2022-28614 \\ \hline
\textbf{CVSS 3.0} & 5.3 \\ \hline
\textbf{Risk Factor} & MEDIUM \\ \hline
\textbf{Description} & The version of Apache httpd installed on the remote host is prior to 2.4.54. Read beyond bounds via \texttt{ap\_rwrite()}: The \texttt{ap\_rwrite()} function in Apache HTTP Server 2.4.53 and earlier may read unintended memory if an attacker can cause the server to reflect very large input using \texttt{ap\_rwrite()} or \texttt{ap\_rputs()}, such as with the \texttt{mod\_lua} \texttt{r:puts()} function. This may lead to information disclosure. \\ \hline
\textbf{Solution} & Upgrade to Apache version 2.4.54 or later. \\ \hline
\end{tabular}

\subsubsection{F-069: Read Beyond Bounds in \texttt{ap\_strcmp\_match()}}
\begin{tabular}{|p{2.8cm}|p{12.2cm}|}
\hline
\textbf{CVE} & CVE-2022-28615 \\ \hline
\textbf{CVSS 3.0} & 9.1 \\ \hline
\textbf{Risk Factor} & CRITICAL \\ \hline
\textbf{Description} & The version of Apache httpd installed on the remote host is prior to 2.4.54. Read beyond bounds in \texttt{ap\_strcmp\_match()}: Apache HTTP Server 2.4.53 and earlier may crash or disclose information due to a read beyond bounds in \texttt{ap\_strcmp\_match()} when provided with an extremely large input buffer. Although no code distributed with the server can directly trigger such a call, third-party modules or Lua scripts using \texttt{ap\_strcmp\_match()} may be affected. \\ \hline
\textbf{Solution} & Upgrade to Apache version 2.4.54 or later. \\ \hline
\end{tabular}

%%%%%%%%%%%%%%%%%%%%%%%%%%%%%%%%%%%%%%%%%%%%
%
% Apache 2.4.x < 2.4.54 Multiple Vulnerabilities
%
%%%%%%%%%%%%%%%%%%%%%%%%%%%%%%%%%%%%%%%%%%%%
\subsection{Apache 2.4.x < 2.4.54 Multiple Vulnerabilities}

\subsubsection{F-070: Denial of Service in mod\_sed}
\begin{tabular}{|p{2.8cm}|p{12.2cm}|}
\hline
\textbf{CVE} & CVE-2022-30522 \\ \hline
\textbf{CVSS 3.0} & 7.5 \\ \hline
\textbf{Risk Factor} & HIGH \\ \hline
\textbf{Description} & The version of Apache httpd installed on the remote host is prior to 2.4.54. Denial of Service in mod\_sed: If Apache HTTP Server 2.4.53 is configured to perform transformations with mod\_sed in contexts where the input to mod\_sed may be very large, mod\_sed may perform excessively large memory allocations and trigger an abort. This issue was reported by Brian Moussalli from the JFrog Security Research team. \\ \hline
\textbf{Solution} & Upgrade to Apache version 2.4.54 or later. \\ \hline
\end{tabular}

%%%%%%%%%%%%%%%%%%%%%%%%%%%%%%%%%%%%%%%%%%%%
%
% Apache 2.4.x < 2.4.54 Multiple Vulnerabilities (mod_lua)
%
%%%%%%%%%%%%%%%%%%%%%%%%%%%%%%%%%%%%%%%%%%%%
\subsection{Apache 2.4.x < 2.4.54 Multiple Vulnerabilities (mod\_lua)}

\subsubsection{F-071: Denial of Service in mod\_lua \texttt{r:parsebody}}
\begin{tabular}{|p{2.8cm}|p{12.2cm}|}
\hline
\textbf{CVE} & CVE-2022-29404 \\ \hline
\textbf{CVSS 3.0} & 7.5 \\ \hline
\textbf{Risk Factor} & HIGH \\ \hline
\textbf{Description} & The version of Apache httpd installed on the remote host is prior to 2.4.54. Denial of service in mod\_lua \texttt{r:parsebody}: In Apache HTTP Server 2.4.53 and earlier, a malicious request to a Lua script that calls \texttt{r:parsebody(0)} may cause a denial of service due to the absence of a default limit on possible input size. This issue was reported by Ronald Crane (Zippenhop LLC). \\ \hline
\textbf{Solution} & Upgrade to Apache version 2.4.54 or later. \\ \hline
\end{tabular}

\subsubsection{F-072: Information Disclosure in mod\_lua with WebSockets}
\begin{tabular}{|p{2.8cm}|p{12.2cm}|}
\hline
\textbf{CVE} & CVE-2022-30556 \\ \hline
\textbf{CVSS 3.0} & 7.5 \\ \hline
\textbf{Risk Factor} & HIGH \\ \hline
\textbf{Description} & The version of Apache httpd installed on the remote host is prior to 2.4.54. Information disclosure in mod\_lua with WebSockets: Apache HTTP Server 2.4.53 and earlier may return lengths to applications calling \texttt{r:wsread()} that point beyond the end of the storage allocated for the buffer, potentially exposing unintended information. \\ \hline
\textbf{Solution} & Upgrade to Apache version 2.4.54 or later. \\ \hline
\end{tabular}

%%%%%%%%%%%%%%%%%%%%%%%%%%%%%%%%%%%%%%%%%%%%
%
% Apache 2.4.x < 2.4.54 Out-Of-Bounds Read (CVE-2022-28330)
%
%%%%%%%%%%%%%%%%%%%%%%%%%%%%%%%%%%%%%%%%%%%%
\subsection{Apache 2.4.x < 2.4.54 Out-Of-Bounds Read (CVE-2022-28330)}

\subsubsection{F-073: Read Beyond Bounds in mod\_isapi}
\begin{tabular}{|p{2.8cm}|p{12.2cm}|}
\hline
\textbf{CVE} & CVE-2022-28330 \\ \hline
\textbf{CVSS 3.0} & 5.3 \\ \hline
\textbf{Risk Factor} & MEDIUM \\ \hline
\textbf{Description} & The version of Apache httpd installed on the remote host is prior to 2.4.54. It is, therefore, affected by an out-of-bounds read vulnerability as referenced in the 2.4.54 advisory. Read beyond bounds in mod\_isapi: Apache HTTP Server 2.4.53 and earlier on Windows may read beyond buffer bounds when configured to process requests with the mod\_isapi module. This issue was reported by Ronald Crane (Zippenhop LLC). \\ \hline
\textbf{Solution} & Upgrade to Apache version 2.4.54 or later. \\ \hline
\end{tabular}

%%%%%%%%%%%%%%%%%%%%%%%%%%%%%%%%%%%%%%%%%%%%
%
% Apache 2.4.x < 2.4.64 Multiple Vulnerabilities
%
%%%%%%%%%%%%%%%%%%%%%%%%%%%%%%%%%%%%%%%%%%%%
\subsection{Apache 2.4.x < 2.4.64 Multiple Vulnerabilities}

\subsubsection{F-074: Core HTTP Response Splitting}
\begin{tabular}{|p{2.8cm}|p{12.2cm}|}
\hline
\textbf{CVE} & CVE-2024-42516 \\ \hline
\textbf{CVSS 3.0} & 7.5 \\ \hline
\textbf{Risk Factor} & HIGH \\ \hline
\textbf{Description} & The version of Apache httpd installed on the remote host is prior to 2.4.64. HTTP response splitting in the core of Apache HTTP Server allows an attacker who can manipulate Content-Type response headers of applications hosted or proxied by the server to split the HTTP response. This issue was originally described as CVE-2023-38709, but the patch included in Apache HTTP Server 2.4.59 did not fully address the vulnerability. The issue is fixed in version 2.4.64. \\ \hline
\textbf{Solution} & Upgrade to Apache version 2.4.64 or later. \\ \hline
\end{tabular}

\subsubsection{F-075: mod\_proxy Server-Side Request Forgery via Content-Type Header}
\begin{tabular}{|p{2.8cm}|p{12.2cm}|}
\hline
\textbf{CVE} & CVE-2024-43204 \\ \hline
\textbf{CVSS 3.0} & 7.5 \\ \hline
\textbf{Risk Factor} & HIGH \\ \hline
\textbf{Description} & The version of Apache httpd installed on the remote host is prior to 2.4.64. Server-Side Request Forgery (SSRF) in Apache HTTP Server with mod\_proxy loaded allows an attacker to send outbound proxy requests to a URL controlled by the attacker. Exploitation requires an uncommon configuration where mod\_headers modifies the Content-Type request or response header using a value provided in the HTTP request. \\ \hline
\textbf{Solution} & Upgrade to Apache version 2.4.64 or later. \\ \hline
\end{tabular}

\subsubsection{F-076: Windows NTLM Hash Leakage via mod\_rewrite / Expressions}
\begin{tabular}{|p{2.8cm}|p{12.2cm}|}
\hline
\textbf{CVE} & CVE-2024-43394 \\ \hline
\textbf{CVSS 3.0} & 7.5 \\ \hline
\textbf{Risk Factor} & HIGH \\ \hline
\textbf{Description} & The version of Apache httpd installed on the remote host is prior to 2.4.64. Server-Side Request Forgery (SSRF) in Apache HTTP Server on Windows allows potential leakage of NTLM hashes to a malicious server through mod\_rewrite or Apache expressions that process unvalidated request input. This issue affects Apache HTTP Server versions from 2.4.0 through 2.4.63. \\ \hline
\textbf{Solution} & Upgrade to Apache version 2.4.64 or later. \\ \hline
\end{tabular}

\subsubsection{F-077: mod\_ssl Log File Escape Character Injection}
\begin{tabular}{|p{2.8cm}|p{12.2cm}|}
\hline
\textbf{CVE} & CVE-2024-47252 \\ \hline
\textbf{CVSS 3.0} & 7.5 \\ \hline
\textbf{Risk Factor} & HIGH \\ \hline
\textbf{Description} & The version of Apache httpd installed on the remote host is prior to 2.4.64. Insufficient escaping of user-supplied data in mod\_ssl in Apache HTTP Server 2.4.63 and earlier allows an untrusted SSL/TLS client to insert escape characters into log files in certain configurations. When CustomLog is configured with \texttt{\%{varname}x} or \texttt{\%{varname}c} to log variables provided by mod\_ssl, such as \texttt{SSL\_TLS\_SNI}, unsanitized client-controlled data may be written to log files. \\ \hline
\textbf{Solution} & Upgrade to Apache version 2.4.64 or later. \\ \hline
\end{tabular}

\subsubsection{F-078: Apache 2.4.x < 2.4.64 Vulnerability (CVE-2025-23048)}
\begin{tabular}{|p{2.8cm}|p{12.2cm}|}
\hline
\textbf{CVE} & CVE-2025-23048 \\ \hline
\textbf{CVSS 3.0} & 9.1 \\ \hline
\textbf{Risk Factor} & CRITICAL \\ \hline
\textbf{Description} & In some mod\_ssl configurations on Apache HTTP Server 2.4.35 through to 2.4.63, an access control bypass by trusted clients is possible using TLS 1.3 session resumption. Configurations are affected when mod\_ssl is configured for multiple virtual hosts, with each restricted to a different set of trusted client certificates (for example with a different SSLCACertificateFile/Path setting). In such a case, a client trusted to access one virtual host may be able to access another virtual host, if SSLStrictSNIVHostCheck is not enabled in either virtual host. \\ \hline
\textbf{Solution} & Upgrade to Apache version 2.4.64 or later. \\ \hline
\end{tabular}

\subsubsection{F-079: mod\_proxy\_http2 Assertion Failure Denial of Service}
\begin{tabular}{|p{2.8cm}|p{12.2cm}|}
\hline
\textbf{CVE} & CVE-2025-49630 \\ \hline
\textbf{CVSS 3.0} & 7.5 \\ \hline
\textbf{Risk Factor} & HIGH \\ \hline
\textbf{Description} & The version of Apache httpd installed on the remote host is prior to 2.4.64. In certain proxy configurations, a denial of service attack against Apache HTTP Server versions 2.4.26 through 2.4.63 can be triggered by untrusted clients causing an assertion failure in mod\_proxy\_http2. Affected configurations require a reverse proxy configured for an HTTP/2 backend with \texttt{ProxyPreserveHost} set to on. \\ \hline
\textbf{Solution} & Upgrade to Apache version 2.4.64 or later. \\ \hline
\end{tabular}

\subsubsection{F-080: Apache 2.4.x < 2.4.64 Vulnerability (CVE-2025-49812)}
\begin{tabular}{|p{2.8cm}|p{12.2cm}|}
\hline
\textbf{CVE} & CVE-2025-49812 \\ \hline
\textbf{CVSS 3.0} & 7.4 \\ \hline
\textbf{Risk Factor} & HIGH \\ \hline
\textbf{Description} & In some mod\_ssl configurations on Apache HTTP Server versions through to 2.4.63, an HTTP desynchronisation attack allows a man-in-the-middle attacker to hijack an HTTP session via a TLS upgrade. Only configurations using "SSLEngine optional" to enable TLS upgrades are affected. \\ \hline
\textbf{Solution} & Upgrade to Apache version 2.4.64 or later. \\ \hline
\end{tabular}

\subsubsection{F-081: Apache 2.4.x < 2.4.64 Vulnerability (CVE-2025-53020)}
\begin{tabular}{|p{2.8cm}|p{12.2cm}|}
\hline
\textbf{CVE} & CVE-2025-53020 \\ \hline
\textbf{CVSS 3.0} & 7.5 \\ \hline
\textbf{Risk Factor} & HIGH \\ \hline
\textbf{Description} & Late Release of Memory after Effective Lifetime vulnerability in Apache HTTP Server. This issue affects Apache HTTP Server: from 2.4.17 up to 2.4.63. \\ \hline
\textbf{Solution} & Upgrade to Apache version 2.4.64 or later. \\ \hline
\end{tabular}

%%%%%%%%%%%%%%%%%%%%%%%%%%%%%%%%%%%%%%%%%%%%
%
% Apache < 2.4.49 Multiple Vulnerabilities
%
%%%%%%%%%%%%%%%%%%%%%%%%%%%%%%%%%%%%%%%%%%%%
\subsection{Apache < 2.4.49 Multiple Vulnerabilities}

\subsubsection{F-082: mod\_proxy Server-Side Request Forgery (SSRF)}
\begin{tabular}{|p{2.8cm}|p{12.2cm}|}
\hline
\textbf{CVE} & CVE-2021-40438 \\ \hline
\textbf{CVSS 3.0} & 9.0 \\ \hline
\textbf{Risk Factor} & CRITICAL \\ \hline
\textbf{Description} & The version of Apache httpd installed on the remote host is prior to 2.4.49. It is, therefore, affected by a vulnerability as referenced in the 2.4.49 changelog. A crafted request URI path can cause mod\_proxy to forward the request to an origin server chosen by the remote user. This vulnerability may allow an attacker to perform Server-Side Request Forgery (SSRF) attacks by controlling the destination of proxied requests. \\ \hline
\textbf{Solution} & Upgrade to Apache version 2.4.49 or later. \\ \hline
\end{tabular}

%%%%%%%%%%%%%%%%%%%%%%%%%%%%%%%%%%%%%%%%%%%%
%
% Apache >= 2.4.17 < 2.4.49 mod_http2
%
%%%%%%%%%%%%%%%%%%%%%%%%%%%%%%%%%%%%%%%%%%%%
\subsection{Apache >= 2.4.17 < 2.4.49 mod\_http2}

\subsubsection{F-083: mod\_http2 Request Splitting and Cache Poisoning}
\begin{tabular}{|p{2.8cm}|p{12.2cm}|}
\hline
\textbf{CVE} & CVE-2021-33193 \\ \hline
\textbf{CVSS 3.0} & 7.5 \\ \hline
\textbf{Risk Factor} & HIGH \\ \hline
\textbf{Description} & The version of Apache httpd installed on the remote host is greater than 2.4.17 and prior to 2.4.49. It is, therefore, affected by a vulnerability as referenced in the 2.4.49 changelog. A crafted method sent through HTTP/2 can bypass validation and be forwarded by mod\_proxy, which can lead to HTTP request splitting or cache poisoning attacks. \\ \hline
\textbf{Solution} & Upgrade to Apache version 2.4.49 or later. \\ \hline
\end{tabular}

%%%%%%%%%%%%%%%%%%%%%%%%%%%%%%%%%%%%%%%%%%%%
%
% Apache >= 2.4.30 < 2.4.49 mod_proxy_uwsgi
%
%%%%%%%%%%%%%%%%%%%%%%%%%%%%%%%%%%%%%%%%%%%%
\subsection{Apache >= 2.4.30 < 2.4.49 mod\_proxy\_uwsgi}

\subsubsection{F-084: mod\_proxy\_uwsgi Out-of-Bounds Read Denial of Service}
\begin{tabular}{|p{2.8cm}|p{12.2cm}|}
\hline
\textbf{CVE} & CVE-2021-36160 \\ \hline
\textbf{CVSS 3.0} & 7.5 \\ \hline
\textbf{Risk Factor} & HIGH \\ \hline
\textbf{Description} & The version of Apache httpd installed on the remote host is greater than 2.4.30 and prior to 2.4.49. It is, therefore, affected by a vulnerability as referenced in the 2.4.49 changelog. A carefully crafted request URI path can cause mod\_proxy\_uwsgi to read beyond the allocated memory boundary, resulting in a server crash and denial of service (DoS). \\ \hline
\textbf{Solution} & Upgrade to Apache version 2.4.49 or later. \\ \hline
\end{tabular}

%%%%%%%%%%%%%%%%%%%%%%%%%%%%%%%%%%%%%%%%%%%%
%
% Apache Default Index Page
%
%%%%%%%%%%%%%%%%%%%%%%%%%%%%%%%%%%%%%%%%%%%%
\subsection{Apache Default Index Page}

\subsubsection{F-085: Apache Default Index Page}
\begin{tabular}{|p{2.8cm}|p{12.2cm}|}
\hline
\textbf{CVE} & N/A \\ \hline
\textbf{CVSS 3.0} & N/A \\ \hline
\textbf{Risk Factor} & N/A \\ \hline
\textbf{Description} & The remote web server uses the default Apache index page. This page may contain sensitive information such as the server root directory and installation paths, which could assist an attacker in gathering information about the server configuration. \\ \hline
\textbf{Solution} & Remove the default index page. \\ \hline
\end{tabular}

%%%%%%%%%%%%%%%%%%%%%%%%%%%%%%%%%%%%%%%%%%%%
%
% Apache ServerTokens Information Disclosure
%
%%%%%%%%%%%%%%%%%%%%%%%%%%%%%%%%%%%%%%%%%%%%
\subsection{Apache ServerTokens Information Disclosure}

\subsubsection{F-086: Apache ServerTokens Information Disclosure}
\begin{tabular}{|p{2.8cm}|p{12.2cm}|}
\hline
\textbf{CVE} & N/A \\ \hline
\textbf{CVSS 3.0} & N/A \\ \hline
\textbf{Risk Factor} & N/A \\ \hline
\textbf{Description} & The HTTP headers sent by the remote web server disclose information that can aid an attacker, such as the server version, operating system, and loaded module versions. This information disclosure may allow attackers to identify the underlying software configuration and target known vulnerabilities associated with the disclosed versions. \\ \hline
\textbf{Solution} & Change the Apache \texttt{ServerTokens} configuration value to \texttt{Prod}. \\ \hline
\end{tabular}

%%%%%%%%%%%%%%%%%%%%%%%%%%%%%%%%%%%%%%%%%%%%
%
% Web Server Error Page Information Disclosure
%
%%%%%%%%%%%%%%%%%%%%%%%%%%%%%%%%%%%%%%%%%%%%
\subsection{Web Server Error Page Information Disclosure}

\subsubsection{F-087: Web Server Error Page Information Disclosure}
\begin{tabular}{|p{2.8cm}|p{12.2cm}|}
\hline
\textbf{CVE} & N/A \\ \hline
\textbf{CVSS 3.0} & N/A \\ \hline
\textbf{Risk Factor} & N/A \\ \hline
\textbf{Description} & The default error page sent by the remote web server discloses information that can aid an attacker, such as the server version and the programming languages or technologies used by the web server. This information disclosure may allow attackers to identify the underlying web server configuration and target known vulnerabilities associated with the disclosed information. \\ \hline
\textbf{Solution} & Modify the web server configuration to avoid disclosing detailed information about the underlying web server, or configure a custom error page instead. \\ \hline
\end{tabular}

%%%%%%%%%%%%%%%%%%%%%%%%%%%%%%%%%%%%%%%%%%%%
%
% Web Server HTTP Header Information Disclosure
%
%%%%%%%%%%%%%%%%%%%%%%%%%%%%%%%%%%%%%%%%%%%%
\subsection{Web Server HTTP Header Information Disclosure}

\subsubsection{F-088: Web Server HTTP Header Information Disclosure}
\begin{tabular}{|p{2.8cm}|p{12.2cm}|}
\hline
\textbf{CVE} & N/A \\ \hline
\textbf{CVSS 3.0} & N/A \\ \hline
\textbf{Risk Factor} & N/A \\ \hline
\textbf{Description} & The HTTP headers sent by the remote web server disclose information that can aid an attacker, such as the server version and the programming languages or technologies used by the web server. This information disclosure may allow attackers to identify the underlying web server configuration and target known vulnerabilities associated with the disclosed information. \\ \hline
\textbf{Solution} & Modify the HTTP headers of the web server to avoid disclosing detailed information about the underlying web server. \\ \hline
\end{tabular}

%%%%%%%%%%%%%%%%%%%%%%%%%%%%%%%%%%%%%%%%%%%%
%
% Linux Kernel Vulnerability
%
%%%%%%%%%%%%%%%%%%%%%%%%%%%%%%%%%%%%%%%%%%%%
\subsection{Linux Kernel Local Privilege Escalation}

\subsubsection{F-089: userns: also map extents in the reverse map to kernel IDs}
\begin{tabular}{|p{2.8cm}|p{12.2cm}|}
\hline
\textbf{CVE} & CVE-2018-18955 \\ \hline
\textbf{CVSS 3.0} & 7.0 \\ \hline
\textbf{Risk Factor} & HIGH \\ \hline
\textbf{Description} & In the Linux kernel 4.15.x through 4.19.x before 4.19.2, map\_write() in kernel/user\_namespace.c allows privilege escalation because it mishandles nested user namespaces with more than 5 UID or GID ranges. A user who has CAP\_SYS\_ADMIN in an affected user namespace can bypass access controls on resources outside the namespace. This occurs because an ID transformation takes place properly for the namespaced-to-kernel direction but not for the kernel-to-namespaced direction. \\ \hline
\textbf{Solution} & Update the Linux kernel to the latest available release. \\ \hline
\end{tabular}